% ════════════════════════════════════════════
% 具体分析（5.4.1）
% ════════════════════════════════════════════
\subsubsection{测算思路}

本部分属于效应测算与情景预测，采用对数尺度预测、Duan 烟雾还原、参数模拟与 Wald 检验相结合的方法。与回归建模关注“模型是否成立”不同，本部分要把系数翻译成运营人员能直接使用的曝光增量、倍数与置信区间。

首先读取回归阶段保存的最终模型 M2，固定一组代表性画像，明确内容类型、主题、月份、星期、时段与粉丝规模，使不同处理组之间只差一个变量。接着在对数尺度上计算线性预测值，再用 Duan 烟雾因子还原到曝光次数，避免对数变换带来的均值低估。然后对 500 条帖子逐条构造付费与自然、视频与图片的反事实，计算样本平均处理效应，并用 1000 次多元正态参数模拟给出置信区间。最后对内容主题做联合 Wald 检验与两两检验，并在六种模型设定下复核付费与视频效应的方向、显著性与量级。具体分析流程如图\ref{fig:analysis4_flow}所示。

\begin{figure}[ht]
    \begin{center}
        \begin{tikzpicture}[x=1cm, y=0.7cm]
            \node[box] (n11) at (0, 0) {读取最终模型};
            \node[box] (n12) at (5, 0) {设定情景画像};
            \node[box] (n13) at (10, 0) {对数尺度预测};
            \node[box] (n22) at (5, -3) {烟雾法还原};
            \node[box] (n21) at (0, -3) {参数模拟};
            \node[box] (n23) at (10, -3) {效应测算};
            \node[box] (n31) at (0, -6) {联合与两两检验};
            \node[box] (n32) at (5, -6) {稳健对比};
            \node[box] (n33) at (10, -6) {结论解释};
            \draw[arrow] (n11) -- (n12);
            \draw[arrow] (n12) -- (n13);
            \draw[arrow] (n23) -- (n22);
            \draw[arrow] (n22) -- (n21);
            \draw[arrow] (n13) -- ++(0,-1.0) -| (n23);
            \draw[arrow] (n21) -- (n31);
            \draw[arrow] (n31) -- (n32);
            \draw[arrow] (n32) -- (n33);
        \end{tikzpicture}
        \caption{效应测算流程图}
        \label{fig:analysis4_flow}
    \end{center}
\end{figure}
